\section{Introduction}

In the 1990s, many forward-looking observers expected that, within 30 years, people would drive autonomous flying cars and live in vertical ultra-modern skyscrapers \cite{Mcrae1994, Berry1996}.
Thirty years later, there is substantial research on autonomous cars \cite{Faisal2019} and sustainable architecture \cite{Sassi2006}, but these forecasts paid less attention to the internet.
The internet allows people to connect with each other anywhere around the world.
In 1994, Jeff Bezos founded Amazon as an online bookshop. Few investors believed in the model, and most people still treated the internet as another bubble.
Amazon can rapidly deliver physical and digital assets to customers through a customer-focused business model. In doing so, it brought online commerce into everyday life \cite{Mellahi2000, Kotha1998}.

A similar shift happened when Satoshi Nakamoto released the Bitcoin whitepaper, which is among the most cited works in the blockchain literature.
Advanced cryptographic methods supported the idea of fully digital, peer-to-peer electronic money. Its supply schedule is fixed and, under the network's current social consensus, unlikely to change, which is why many people refer to Bitcoin as digital gold \cite{Taskinsoy2021, Uddin2020}. Nevertheless, scepticism remains \cite{Nakamoto}.

The internet can be read as having three asset-related phases.
The first phase placed real-world assets, such as cars or books, into the virtual world.
The second phase enhanced these assets with features that are possible only in digital environments.
Using augmented reality, users can place virtual furniture from IKEA in their homes without visiting the store or guessing how it would look in the room.
The third and latest phase is the emergence of digital-first assets. The first internet-native assets in this sense are cryptocurrencies. Unlike traditional assets that derive value from physical scarcity or institutional backing, cryptocurrencies rely on cryptographic guarantees, fixed supply schedules, and decentralized consensus for their value proposition. This fundamental shift challenges conventional notions of money: without a central authority guaranteeing redemption and with high price volatility, adoption requires trust in mathematical proofs rather than institutional promises \cite{Taskinsoy2021}.

Moving towards an internet-native world, humans are not the only entities creating and sharing value in virtual space. With advances in large language models, AI agents can perform parts of the same behaviour. This creates a need for digital environments in which humans and AI agents can co-exist. The common term for such virtual worlds is the Metaverse, first used in the science-fiction novel \emph{Snow Crash} in 1992 \cite{Stephenson1992}.
At the time of writing, the author predicted the use of electronic currency that was fast, secure, and easy to use. In contrast to traditional "hyperinflated" currencies, digital currency has become a better monetary asset for storing and exchanging value.
Nevertheless, why is digital currency so valuable? The answer cannot be condensed into a single statement.
The most significant advantage of electronic currency is the ability to build a set of components that can later be reused and composed into other parts of the system.
Why mention \emph{Snow Crash} in a thesis on the Internet of Value? The novel imagined a world where information, reputation, and access function as currency, which is close to what Internet of Value infrastructure now makes technically possible. In that sense, the novel already contained an early use case for the Internet of Value: a marketplace for buying and selling information \cite{Cao2022, Kraus2022}.
Non-fungible data can carry different kinds of value and can be exchanged for tokens with monetary value without traditional intermediaries.

% para economic
% trading with information
% virtual reality
% virtual worlds
% language processing
% meta language
% meta
% https://en.wikipedia.org/wiki/Snow_Crash

However, while centralized databases can technically store and transfer digital assets, they cannot provide the decentralized trust guarantees required for a truly open, peer-to-peer value network. A centralized server remains under the control of a single operator who can unilaterally modify records, censor transactions, or become a single point of failure \cite{Chowdhury2018}.
Using blockchain as a source of truth, we can create a decentralized system where data is mutable only through agreement among network participants \cite{Zheng2017}.
Every modification is recorded in an append-only list of transactions that can be read backwards, which makes audit and analysis more straightforward \cite{Wang2022}.

The Internet of Value is no longer a thesis-stage idea. Treasury systems such as Polkadot OpenGov already process hundreds of funding proposals each year \cite{Valastin2024proto}. Decentralised finance protocols hold tens of billions of dollars across heterogeneous tokens \cite{Schar2021}. Tokenised assets now span fungible currencies, non-fungible collectibles, governance rights, and real-world claims \cite{Bambridge2022}. The cost of misjudging value in this setting is concrete. Voters approve grants without reading the underlying evidence. Users hand signing authority to scripts whose behaviour they cannot inspect, and a single misconfigured contract call is enough to drain their account. Communities undervalue assets whose worth lies in utility or contribution rather than in floor price. Large language models and tool-orchestration standards now make it technically feasible to delegate parts of this work to autonomous agents \cite{Sumers2024, Wu2023}, but the question this thesis takes up is not whether such delegation is possible. It is on what terms it is defensible. Agents that act on-chain on behalf of users must interpret value under a stated intent, expose the evidence behind that interpretation, and operate inside a policy that humans can read, audit, and override. Without an auditable value-interpretation layer and a policy boundary that gates authorisation on it, agentic action on-chain reproduces the very intermediaries that the Internet of Value set out to remove.

The research problem we address sits at the intersection of three gaps that current platforms and the surrounding literature do not close jointly. First, asset valuation in the Internet of Value is reduced to market price, even though the value of many on-chain assets lies in utility, governance weight, access rights, contribution to a public good, or productive yield \cite{Lohmer2024, Ante2021}. Second, existing account models, including recent account-abstraction proposals such as EIP-4337 \cite{Buterin2021, Wang2023, Duan2024}, are built around flexible authentication and gas sponsorship rather than around account-level enforcement of audit-trace-gated, policy-constrained agent action. Third, on-chain governance assumes that human voters can evaluate proposals at the rate at which they arrive, an assumption that breaks once treasury volume scales \cite{Valastin2024proto}. We argue that these three gaps are not independent and cannot be closed by adding another wallet, another oracle, or another voting interface. They share a common cause, the absence of a formal decision layer that turns stated intent into auditable, evidence-grounded value interpretation before any value-relevant action is authorised on-chain.

The central contribution of this thesis is Intent-Conditioned Value Projection (ICVP), a formal decision model in which a stated user or governance intent acts as a projection over heterogeneous Internet of Value evidence spaces. Under ICVP, intent does not only change how much Availability, Durability, and Utility matter. It changes what counts as Availability, Durability, and Utility in the first place. The same Pudgy Penguin token valued for speculative trading, for community access, and for use as DeFi collateral pulls in different evidence and instantiates the three valuation dimensions differently in each case. The same is true for stETH evaluated for long-term holding, for collateralisation, or for liquidity exit, and for a Polkadot OpenGov proposal evaluated under public-resource stewardship rather than private acquisition. ICVP makes this projection explicit, produces an auditable valuation trace that records the stated intent, the selected evidence with its provenance, the intent-instantiated dimensions, and the decision in $\{\textit{proceed}, \textit{defer}, \textit{decline}\}$, and exposes the trace to deterministic validation before a human approves the corresponding action.

The Internet of Value makes assets transferable. ICVP, in this thesis, is the layer at which their value becomes interpretable under intent. The agentic system proposed here puts ICVP into practice, while the Non-Fungible Token as Account (NFTAA) primitive proposed in this work acts as the account-level enforcement layer. NFTAA does not compute value. It enforces that only human-approved, policy-compliant actions linked to a valid valuation trace can affect assets on-chain. The layered agentic architecture in Chapter~\ref{sec:design} is the vehicle through which these contributions are tested, not the contribution in itself \cite{Xi2023}. Chapter~\ref{sec:evaluation} reports on whether ICVP improves decision quality, evidence relevance, and auditability across NFT, productive-asset, and treasury-governance scenarios relative to single-dimensional and fixed-weight valuation baselines.

This work is organized as follows. Chapter \ref{sec:blockchain} describes the fundamentals and advanced concepts of blockchain technology.
Chapter \ref{sec:analysis} explains why the Internet of Value matters and identifies the gaps in current platforms that motivate this work.
The detailed aims of this work are described in Chapter \ref{sec:aims}.
The proposed architecture and its response to the identified gaps are presented in Chapter \ref{sec:design}, with a comparison to related work.
Finally, Chapter \ref{sec:conclusions} discusses the results, conclusions, and future work.
